% alt: Three wires for X_{j-1}, the ring X_j minus X_{j-1}, and the outside Lambda minus X_j; the inverse filter on the inner two wires, then the split term h_Y on the outer two wires, then the filter on the inner two wires, with a line below saying which filters of K were stripped and which commuted away.
% Section 4 (part1_arealaw.tex, proof of Lemma lem:buffer, paragraph "Energy of the optimized vector"): for Y split by X_j the conjugated term <phi, K h_Y K^{-1} phi> with K = L_m ... L_1 reduces to <phi, L_j h_Y L_j^{-1} phi>, because the outer filters L_{j+1}, ..., L_m are stripped from the outside in by the stripping identity eq:strip and the inner filters L_1, ..., L_{j-1} commute with h_Y (Lemma lem:singlesplit: Y is contained in X_i for i > j and disjoint from X_i for i < j).
% Ink-to-index: the top wire is X_{j-1}, the middle wire the ring X_j \ X_{j-1}, the bottom wire Lambda \ X_j; L_j^{-1} and L_j are the two-wire gates on the top two wires (their support X_j), h_Y the two-wire gate on the bottom two wires (Y meets the ring and the outside but not X_{j-1}). Read left to right: L_j^{-1} acts first, then h_Y, then L_j; all wires are contracted with phi on both ends.
\[
  \begin{gathered}
    \langle\phi,Kh_YK^{-1}\phi\rangle
    \;=\;
    \Bigl\langle\phi,\;
    \begin{quantikz}
      \lstick{$X_{j-1}$} & \gate[2]{L_j^{-1}} & & \gate[2]{L_j} & \\
      \lstick{$X_j\setminus X_{j-1}$} & & \gate[2]{h_Y} & & \\
      \lstick{$\Lambda\setminus X_j$} & & & &
    \end{quantikz}
    \;\phi\Bigr\rangle
    \\[2pt]
    \text{\small outer filters $L_{j+1},\dots,L_m$: stripped by the stripping identity, $[L_k,\rho_k]=0$}
    \\[-1pt]
    \text{\small inner filters $L_1,\dots,L_{j-1}$: disjoint from $Y$, commute with $h_Y$}
  \end{gathered}
\]
